\section{The most expensive mistakes}
\label{sec:mistakes}

Each of them has cost somebody an evening. They are listed in the order they tend to happen.

\begin{booktable}{@{}rLL@{}}
\thead{\#} & \thead{Mistake} & \thead{Symptom} \\ \midrule
1 & The parser does not check the length field against the size of the buffer & Buffer overflow;
  all sorts of other things break \\
2 & The parser goes back to \code{WaitForSof1} on a second \code{0xA5} & The odd frame is missed,
  inexplicably \\
3 & Multi-byte fields are serialised by copying memory & Works until the other side is built with
  a different endianness \\
4 & A duplicate is not acknowledged again & The sender retransmits for all eternity \\
5 & \code{TX_DATA_LO} and \code{TX_DATA_HI} swapped & The frame arrives --- back to front \\
6 & The cast after the shift in \code{readReg()} & Works for small values, breaks at bit 7 \\
7 & \code{TX_SEND} written before the data registers & The node runs exactly one frame behind \\
8 & \code{clearError()} writes \code{0x1} & The error bit never goes low; all error handling stops
  working \\
9 & \code{RX_ACK} forgotten & The same frame is read in an endless loop \\
10 & \code{receive()} does not mask against \code{dlc} & Up to six bytes are reported that were
  never sent \\
11 & \code{while (!send(frame)) {}} on an invalid frame & Hangs forever \\
12 & The pin mux not set & The program runs, the wires are silent, no errors are reported \\
13 & The data register is not read in \code{transfer()} & All the data is one byte out \\
14 & \code{VDDIO2} not connected & Looks exactly like a broken SPI implementation \\
15 & The factory's members in the wrong declaration order & A reference to an object that has not
  been constructed \\
\end{booktable}

\begin{aside}
\note Eleven of the fifteen would have been caught by a test on the laptop. Three of them --- 12,
13 and 14 --- would not, and those are precisely the three that live in the one class no automated
test covers. That is no coincidence, and it is the whole argument for the ladder in
\kapref{ch:bringup}.
\end{aside}

\section{After the course}

Two things continue directly from here.

\textbf{Testing.} Throughout the project you have run and read tests somebody else wrote. The next
step is deciding what they should check --- which is a harder problem than it sounds, and the one
that decides whether a test suite is worth anything.

\textbf{Your own repository.} The driver you have built is a complete, tested, layered driver
stack against real hardware, developed against a contract and verified without access to the other
half. That is a reasonable first item in a portfolio, and a more realistic description of embedded
development than most course projects.
